\section{Implementación}\label{sec:implementacion-laura}
\subsection{Preparación y compatibilidad del esquema}
\texttt{database/schema.sql} crea las siete tablas con claves primarias autoincrementales, claves foráneas, valores predeterminados y restricciones. \texttt{db:setup} prepara una base nueva y comprueba su compatibilidad. \texttt{db:seed} crea el administrador y paquetes de demostración; \texttt{db:demo} prepara un cliente y una reserva de una clase terminada para ejercitar la asistencia.

El arranque comprueba la conexión antes de escuchar peticiones. Se amplió \texttt{checkSchema} para detectar columnas faltantes, identificadores sin autoincremento y columnas heredadas obligatorias sin valor predeterminado. Evita anunciar una API lista cuando las inserciones fallarán por esas diferencias. No modifica ni migra datos de una base existente, y no certifica todos los tipos o restricciones; una migración requiere comparar el esquema completo.

\subsection{Compra y protección del precio}
En \texttt{purchase-service.js}, el cuerpo solo admite \texttt{packageId}. El usuario comprador se obtiene de la sesión y el precio del paquete se lee dentro de la transacción. El fragmento siguiente pertenece al servicio real; no permite sustituir el precio con un valor enviado por el comprador.
\begin{lstlisting}[caption={Precio de compra elegido por el servidor.}]
return transaction(async (connection) => {
  const item = found(await packages.findById(
    packageId, connection, true), 'Package');
  return repository.create(user.id, packageId,
    item.price, connection);
});
\end{lstlisting}

\subsection{Reserva, asistencia y reseña}
\texttt{reservation-service.js} verifica horario, duplicidad y cupo antes de insertar. El bloqueo de la clase mantiene la decisión de cupo dentro de la transacción. El siguiente fragmento aplica el rechazo cuando no quedan plazas.
\begin{lstlisting}[caption={Control de cupo en la transacción de reserva.}]
const total = await classes.reservationCount(
  classId, connection);
if (total >= item.capacity) {
  fail(409, 'The class is full');
}
return repository.create(user.id,
  classId, connection);
\end{lstlisting}
El administrador registra \texttt{attended} o \texttt{no\_show} después de terminar la clase. \texttt{review-service.js} exige la reserva del usuario con estado \texttt{attended}; rechaza una clase no asistida y una reseña duplicada. La calificación es obligatoria. Comentario e intensidad pueden omitirse; las etiquetas omitidas se convierten en un arreglo vacío. Eliminar una reseña conserva la reserva y la asistencia.

\subsection{Progreso y marcas personales}
\texttt{progress-service.js} usa únicamente reservas asistidas, agrupa por categoría y deduplica semanas para calcular la racha. La semana inicia el lunes y la fecha de negocio usa \texttt{America/Bogota}; una clase se asigna a la semana de su inicio. La racha puede continuar desde la semana anterior si la semana actual aún no tiene asistencia.

\texttt{record-service.js} conserva ejercicio, tipo, valor, unidad y fecha. Rechaza valores no positivos y fechas futuras; el propietario puede crear, modificar o eliminar su marca. La API entrega el historial registrado, sin atribuirle un cálculo automático de máximos ni comparaciones entre unidades diferentes.

\subsection{Lectura del código y alcance}
Los fragmentos muestran las decisiones que conectan el modelo con las reglas. La ejecución automática, las mediciones y las respuestas capturadas pertenecen a la sección de pruebas del informe general. El backend permite construir las vistas del cliente y del administrador, pero la existencia de los endpoints no equivale a una interfaz implementada.
